\scanpage{40}
$U_\rho$ opère (\unsure{sans unité}) sur $D_0$ et $D_0^{**}$,
et que $D_0^{**}$ est un \uline{torseur} sous $U_\rho$
(trivialisé par $\unsure{b_{0\ldots n}}$).

\figure{\begin{tikzpicture}\draw (0,0)--(2,0);\end{tikzpicture}}

Supposons maintenant les $\mu_i=\unsure{2}$ de sorte qu'on a une forme quadratique
\uline{inv.} canonique $f$ sur $\mathcal E$, unique par la condition
\[
f(a_{\unsure{-1}})=0,\qquad f(e_0)=1,
\]
\deleted{\unsure{Je ai}}
d'où $Q\subset X=\mathbb P^{\unsure{\vee}}(\mathcal E)$.

On sait que la \unsure{ps. réfle}
$\operatorname{id}+\underbrace{\unsure{\lambda}\Pi'\otimes\Pi}_{\Pi'\otimes\lambda\Pi}$
est \deleted{\unsure{une}} orthogonale sss l'opération \unsure{corr.}
dans $X$ conserve $Q$, et ceci signifie que
\[
\lambda\,\deleted{\unsure{\Psi}}\inserted{f}(\lambda\Pi)\Pi'
+\Psi(\lambda\Pi)=0\quad\text{i.e.}
\]
\[
\lambda\left[
\lambda\underbrace{f(\Pi)}_{\inserted{-\widetilde\sigma\widetilde\rho}}\Pi'
+\underbrace{\Psi(\Pi)}_{\inserted{-\sigma\Pi'}}
\right]=0\qquad\text{i.e.}
\]
\deleted{$\unsure{\lambda\rho\Pi'+\Psi}$}
\[
\unsure{(\alpha b)}\qquad
\lambda(\sigma+\lambda\widetilde\sigma\widetilde\rho)=0
\]
\mspage{39}

